% Convex combination (a segment) and convex hull (the rubber-band set).
%
%   figures/tikz/ra-convex-combination.tex
%       -> media/figures/ra-convex-combination.{png,svg}
%
% Lecture 13 (Real Analysis Review), convexity. Revisits the convex hull used in
% Lectures 4 and 5 (logical modelling, integer programming) with a definition.
%
% GEOMETRY (1 unit = 1 cm).
%  Left: x_a = (0.3,0.5), x_b = (3.0,1.7). The point theta x_a + (1-theta) x_b
%   is marked for theta = 1, 3/4, 1/2, 1/4, 0:
%     theta = 3/4 -> (0.975,0.800);  1/2 -> (1.650,1.100);
%     1/4 -> (2.325,1.400);          0 -> (3.000,1.700).
%  Right: S = {(0.3,0.3),(2.8,0.1),(3.4,1.4),(2.2,2.7),(0.5,2.2)} are the
%   vertices of the hull (convex pentagon, checked by cross products);
%   (1.5,1.2),(2.0,0.7),(1.0,1.7) are strictly inside, so they add nothing to
%   the hull.
%
% GREYSCALE: points differ by filled vs open marker and by label; the hull is
% hatched AND bounded by a solid line.
\usetikzlibrary{patterns}
\definecolor{okabeblue}{HTML}{0072B2}

\begin{tikzpicture}[
    >={Latex[length=2.0mm]},
    font=\small,
    pt/.style={fill=black, draw=black, circle, inner sep=0pt, minimum size=3.6pt},
    hpt/.style={fill=white, draw=black, circle, inner sep=0pt, minimum size=3.6pt},
    lab/.style={font=\scriptsize, inner sep=1.5pt},
    panelcap/.style={font=\small, align=center, text width=52mm, anchor=north},
  ]
% ---------------------------------------------------------- left: segment
\begin{scope}
  \draw[okabeblue, line width=1.3pt] (0.3,0.5) -- (3.0,1.7);
  \node[pt] (xa) at (0.3,0.5) {};
  \node[pt] (xb) at (3.0,1.7) {};
  \node[lab, below right] at (xa) {$\mathbf{x}_a$};
  \node[lab, above left] at (xb) {$\mathbf{x}_b$};
  \node[lab, above left] at (xa) {$\theta = 1$};
  \node[lab, below right] at (xb) {$\theta = 0$};
  \foreach \th/\px/\py in {0.75/0.975/0.800, 0.5/1.650/1.100, 0.25/2.325/1.400} {
    \node[hpt] at (\px,\py) {};
  }
  \node[lab, above left] at (0.975,0.800) {$\tfrac34$};
  \node[lab, above left] at (1.650,1.100) {$\tfrac12$};
  \node[lab, above left] at (2.325,1.400) {$\tfrac14$};
  \node[panelcap] at (1.65,-0.15) {\textbf{convex combinations}\\[-1pt]
        $\theta\mathbf{x}_a + (1-\theta)\mathbf{x}_b$, $0 \le \theta \le 1$,\\[-1pt]
        fill the segment};
\end{scope}
% ---------------------------------------------------------- right: hull
\begin{scope}[xshift=5.4cm]
  \path[pattern=north east lines, pattern color=okabeblue!55]
     (0.3,0.3) -- (2.8,0.1) -- (3.4,1.4) -- (2.2,2.7) -- (0.5,2.2) -- cycle;
  \draw[okabeblue, line width=1.2pt]
     (0.3,0.3) -- (2.8,0.1) -- (3.4,1.4) -- (2.2,2.7) -- (0.5,2.2) -- cycle;
  \foreach \px/\py in {0.3/0.3, 2.8/0.1, 3.4/1.4, 2.2/2.7, 0.5/2.2} {
    \node[pt] at (\px,\py) {};
  }
  \foreach \px/\py in {1.5/1.2, 2.0/0.7, 1.0/1.7} {
    \node[hpt] at (\px,\py) {};
  }
  \node[lab, anchor=west, fill=white, inner sep=1pt] at (2.15,1.55) {$\operatorname{conv}(S)$};
  \node[panelcap] at (1.85,-0.15) {\textbf{convex hull} $\operatorname{conv}(S)$\\[-1pt]
        the smallest convex set\\[-1pt] containing $S$};
\end{scope}
\end{tikzpicture}
